% Generated draft: requires mathematical and source review.
\documentclass[11pt]{article}
\usepackage[margin=25mm]{geometry}
\usepackage{fontspec}
\setmainfont{Latin Modern Roman}
\usepackage{amsmath,amssymb,mathtools,unicode-math}
\setmathfont{Latin Modern Math}
\usepackage{enumitem,needspace,xurl,hyperref}
\setlength{\parindent}{0pt}
\setlength{\parskip}{5pt}
\setlength{\emergencystretch}{3em}
\newcommand{\N}{\mathbb{N}}
\newcommand{\Z}{\mathbb{Z}}
\newcommand{\Q}{\mathbb{Q}}
\newcommand{\R}{\mathbb{R}}
\newcommand{\C}{\mathbb{C}}
\newcommand{\F}{\mathbb{F}}
\newcommand{\A}{\mathbb{A}}
\newcommand{\PP}{\mathbb{P}}
\newcommand{\E}{\mathbb{E}}
\newcommand{\eps}{\varepsilon}
\newcommand{\dd}{\,\mathrm{d}}
\begin{document}
\section*{UnsolvedMath ID: 30005030}
\textbf{Source problem code:} catalogue entry
\subsection{Definitions and mathematical statement}
Let \(W^H\) denote fractional Brownian motion with Hurst parameter \(H\in(0,1)\). In one real component, this is a centered Gaussian process with covariance
\[
\mathbb E[W_t^H W_s^H]=\tfrac12\bigl(t^{2H}+s^{2H}-|t-s|^{2H}\bigr).
\]
The supplied record does not specify the spatial dimension, spatial domain, time interval, initial condition, or any narrower admissible range of \(H\).

The notation \(b\in L_t^q C_x^\alpha\) means that the time-dependent drift belongs to an \(L^q\) space in time with values in a spatial regularity space denoted by \(C^\alpha\). The precise convention for \(C^\alpha\), including its meaning for nonpositive or integer \(\alpha\), and the admissible range of \(q\) are not supplied.

The extracted question concerns the formal stochastic differential equation
\[
\dd X_t=b(t,X_t)\,\dd t+\dd W_t^H.
\]
Can increasing the time integrability exponent \(q\) in the assumption \(b\in L_t^q C_x^\alpha\) permit a lower spatial regularity threshold than
\[
\alpha>1-\frac{1}{2H},
\]
in accordance with the source's scaling prediction?

This is an unresolved formulation of the supplied question, rather than a fully specified theorem target. The source excerpt refers to formula (2), but that formula is absent from the supplied record, so the predicted relation among \(\alpha,q,H\) cannot be stated from the available evidence. The record also does not identify the solution concept or the property to be established, such as existence, pathwise uniqueness, or another regularization property. If the intended spatial regularity allows distributional drifts, the meaning of the drift term requires additional specification. No particular interpretation of these missing qualifications is imposed here.
\subsection{Short English statement}
The supplied question asks whether stronger time integrability of the drift can compensate for weaker spatial regularity in an equation driven by fractional Brownian motion, allowing spatial exponents below the stated threshold. The scaling formula, solution concept, parameter ranges, and precise mathematical conclusion are missing, so a complete quantified statement requires review of the source context.
\subsection{Sources}
\begin{itemize}
\item Local UnsolvedMath registry, supplied source record with permanent ID \(30005030\), title ``Spatial Regularity Versus Time Integrability for Fractional SDEs,'' source problem code \texttt{catalogue entry}, in the Oberwolfach Reports Open Problems collection. No individual proposer is supplied; the proposed year is recorded as 2022.
\item Supplied report citation: \emph{Mini-Workshop: Regularization by Noise: Theoretical Foundations, Numerical Methods and Applications} (2023), Oberwolfach Reports, pp.~457--488, DOI 10.4171/owr/2022/9. Supplied source URL and literature source labeled ``Oberwolfach report'': \url{https://doi.org/10.4171/owr/2022/9}. The supplied citation's year 2023 and the record's proposed year 2022 are retained without reconciliation.
\item Literature assessment dated \texttt{2026-08-22}: the supplied curation labels the problem open and reports that the source leaves the time-integrability improvement question open, with no later result located settling this exact threshold question. This is supplied dated curation, not an independently verified literature assessment.
\item The registry labels the statement \texttt{corrected\_verified} and reports that the official report was checked directly when preparing the clean statement. That verification claim is part of the supplied record; no independent source visit or verification is asserted here. The retained original excerpt contains extraction artifacts and refers to an unavailable formula (2).
\end{itemize}
\end{document}
